\scanpage{21}
[P. ex. dans le cas où $c$ existe \inserted{\text{-- qu'on prendra pour origine --}}, $G$ est l'ens des homothéties de centre $c$ (isomorphe à $k^*$) qui est évidemment simplement transitif sur l'ens des pts de la droite \unsure{homogène} \uline{épointée}.]
% slanted
\margin{c'est le cas \deleted{\unsure{ssi}} $n\ne\operatorname{car} k$ \deleted{\unsure{$n\ne p$}} et $n\ne4$ si $p=2$ (ou $6$)}
\inserted{\text{En fait il est clair \inserted{\text{a priori}} que $G$ opère librement, à cause \deleted{du} cor \unsure{?} \deleted{\unsure{(prop)}} de la prop. 1. Donc \deleted{\unsure{le montrer plus court}} la question est celle de la transitivité\ldots}}
% slanted
\margin{on trouve déjà que $k^*=G$ dans le cas où $c$ ``existe'' (i.e. centre pas à l'infini!) Dans le cas $D_\sigma\parallel D_{\sigma'}$, on \deleted{\unsure{trouvait}} translations de direction $D_\sigma$ commutent à $\deleted{\unsure{D}}\sigma,\sigma'$ d'où la conclusion}

\uline{Cor 1}) Deux $n$-polygônes réguliers \deleted{\unsure{g}} \unsure{épinglé|épinglés}, tels que les représentations correspondantes de $D_n$ soient isomorphes, sont isomorphes. (la réciproque étant triviale). Donc la classific\unsure{ion|ation} mod isomorphismes des représentations affines \inserted{\text{admissibles}} de dim $2$ de $D_n$ (sur $k$ corps fixé) équivaut à celle des $n$-polygônes réguliers.

\uline{Cor} $2)$ Soit
\[
\begin{aligned}
\varphi &:D_n\mathbin{\unsure{-|\to}}\mathrm{Aff}(E)\\
\varphi' &:D_n\mathbin{-}\mathrm{Aff}(E')
\end{aligned}
\]
deux repr. admissibles, et soit $P$ un polygone \unsure{$n$-régulier|$n$régulier} subordonné à $\varphi$ (dont la donnée équivaut à celle d'un $s_0\in D_\sigma^*\ldots$). Alors si à tt isom $\Theta:(E,\varphi)\to(E',\varphi')$ (isom. affine \deleted{\unsure{deux}} tel que $\Theta\varphi(g)\Theta^{-1}=\varphi'(g)$ $\forall g\in D_n$) on associe le polygône régulier \deleted{\unsure{sub}} $P'=\Theta(P)$, subordonné à $P'$, on
% archive sheet number: 20
